(1) 相邻狭缝的远场光程差为 $d\sin\theta$，相位差为 $\delta=2\pi d\sin\theta/\lambda$。主极大要求 $\delta=2\pi m$，即 $d\sin\theta_m=m\lambda$，其中 $m\in\mathbb Z$ 且 $|m\lambda|\le d$ 才有传播级次。
(2) 固定整数级次 $m$，对 $d\sin\theta_m=m\lambda$ 微分得 $d\cos\theta_m\,d\theta_m=m\,d\lambda$，故 $d\theta_m/d\lambda=m/(d\cos\theta_m)$。在非零级次和衍射角较大时角分散率提高，但 $\cos\theta$ 趋零时近似仪器设计需注意边缘。
(3) $N$ 缝衍射强度因子为 $[\sin(N\delta/2)/\sin(\delta/2)]^2$，主极大周围的第一个零点与主极大相位差约 $2\pi/N$。Rayleigh判据要求两近邻波长的同级主峰角差等于一条主峰到首零点的角差，得到 $\Delta\lambda=\lambda/(|m|N)$，分辨本领 $\lambda/\Delta\lambda=|m|N$。
(4) 零级 $m=0$ 不随波长偏转，因此没有有效的角色散；在非零高阶近掠射角时有限孔径与几何畸变影响显著。增大实际被照亮刻线数 $N$ 会使主峰变窄并改善分辨，但若只增大未照明的刻线数，不会得到公式预测的提升。
